\subsection{Watershed Segmentation for Crowded-Ball Detection}
\label{sec:watershed-experiment}

The first segmentation-based detection experiment investigated whether marker-controlled watershed
segmentation could recover balls missed by the circular Hough detector in
crowded regions. Unlike the preceding experiments, its primary objective
was improved localisation rather than a different type or colour threshold.
The method was implemented as an independent fixed prototype in
\code{experiments/watershedDetection.cpp}, without replacing the detector
used by the main application. No parameter grid was run.

Watershed separates image regions by growing regions from labelled markers.
An additional ball therefore requires an additional suitable marker; merely
running watershed from the existing Hough centres cannot recover a ball
whose interior has no marker. The experiment followed the general
distance-transform/marker approach illustrated in
\href{https://docs.opencv.org/4.12.0/d2/dbd/tutorial_distance_transform.html}
{OpenCV's distance-transform and watershed tutorial}, adapted to the known
table geometry and existing Hough detections.

\subsubsection{Auditing the baseline misses}

The current detector was rerun on every image, including table masking,
Gaussian smoothing, Hough detection and the existing pocket filter. Its
centres were transformed into table coordinates and matched one-to-one
to annotations within 15 table pixels. Raw Hough candidates were evaluated
separately to identify whether the pocket filter removed a previously
matched ball. The audit also recorded whether a missed annotation lay
inside one of the image-derived crowded-region windows.

% Generated by experiments/writeRecognitionReport.py from saved CSVs.
\begin{table}[htbp]
\centering
\small
\caption{Baseline missed annotations. Source coordinates are obtained by inverse perspective transformation for evaluation and visualisation only; they never supply watershed markers.}
\label{tab:watershed-misses}
\begin{tabular}{lrrrrcc}
\toprule
Image & Ball & Source x & Source y & Nearest error & Raw match & In ROI \\
\midrule
Table1.jpeg & 6 & 1403.6 & 149.7 & 33.47 & No & Yes \\
Table1.jpeg & 7 & 1431.9 & 202.7 & 33.74 & No & Yes \\
Table3.jpeg & 10 & 1474.9 & 101.2 & 55.43 & No & Yes \\
Table4.jpeg & 10 & 1474.9 & 101.2 & 52.49 & No & Yes \\
\bottomrule
\end{tabular}
\end{table}
\FloatBarrier


On Table1--Table6, the current run matched 86 of 90 annotations and produced
two false detections. The four misses were balls 6 and 7 in Table1 and
ball 10 in both Table3 and Table4. All four were already unmatched before
pocket filtering and all lay inside crowded windows. Thus this run did
not support pocket filtering as the immediate explanation for these misses.
The annotations' distances to their nearest filtered detections exceeded
the 15-pixel matching tolerance, rather than simply requiring a different
assignment of two otherwise eligible matches.

These counts were measured by rerunning the current source and agree with
the current baseline in Table~\ref{tab:detection-results}. An older saved
validation CSV recorded seven false detections; it is retained as a historical
artifact rather than mixed with these new comparison results.
Recognition parameters had already been
tuned using Table1--Table6; downstream recognition scores therefore remain
exploratory rather than independent test estimates.

Inverse perspective transformation was used to display missed annotations
at their approximate original-image positions. These points were used
only for evaluation and diagnostic overlays. They did not influence
window construction, foreground masks, distance peaks, markers, basin
validation or proposal acceptance. The segmentation function accepts an
image and baseline circles, not annotations.

\subsubsection{Local windows and foreground construction}

Pairs of filtered Hough circles were considered nearby when their centre
distance was at most 1.4 times the sum of their radii. Transitive groups
with at least two detections formed the initial crowded groups. Their
circle bounding boxes were expanded by twice the median group radius
and clipped to the image. Overlapping windows were merged, and every
baseline circle centred in a resulting window supplied a Hough marker.
The expected radius for new candidates was the median Hough radius in
that window. Ten windows were processed across the seven images.

This local design preserves isolated detections and limits processing of
unrelated cloth or table features. Its limitation is equally explicit:
a completely missed group with no suitable nearby Hough detections, or
an isolated missed ball outside the selected windows, cannot be recovered.
The audit showed that this coverage limitation did not exclude any of
the four missed annotations in this particular run.

Cloth colour was estimated separately in each window. Eligible samples
lay inside the known table polygon, outside disks of 1.25 times each
local Hough radius, and satisfied $35\leq H\leq95$, $S\geq35$ and
$V\geq30$ in OpenCV HSV units. Samples were taken every two pixels in
each direction. At least 40 eligible samples were required; otherwise
the window was skipped without changing the Hough detections.

The samples were converted to OpenCV's encoded 8-bit Lab representation.
The cloth representative was the component-wise median Lab vector. Let
$d_j$ be a sample's Euclidean distance from that representative. A pixel
was provisionally foreground when its Lab distance exceeded
\[
 T=\operatorname{clip}_{[16,45]}\left(
      \operatorname{median}(d_j)
      +2.5\operatorname{median}\left|d_j-\operatorname{median}(d_j)\right|
    \right).
\]
These numerical distances apply to the encoded channel values and should
not be interpreted as calibrated perceptual colour differences.

Small foreground cores of radius $0.65r$ were retained around existing
Hough centres. This reduced the chance that a known green ball disappeared
entirely into a green-cloth mask. It did not solve the corresponding
problem for a missed green ball with no Hough marker. A $3\times3$
elliptical closing operation filled narrow gaps, enclosed holes were
filled using flood filling from a padded exterior, and the result was
intersected with the table polygon. In a tightly packed group, this
foreground can remain one merged object rather than one object per ball.

\subsubsection{Markers, watershed and candidate acceptance}

Three variants were compared:
\begin{enumerate}
\item \textbf{Hough baseline:} the unchanged filtered Hough circles.
\item \textbf{Hough-seeded watershed:} one marker per existing Hough
detection, with no new detection proposals. This is a segmentation
control: circles are not relocated and no Hough circle is removed.
\item \textbf{Hybrid watershed:} Hough markers plus additional
distance-transform peaks that pass placement and basin checks.
\end{enumerate}

The foreground's Euclidean distance transform measured distance to the
nearest zero pixel. Local maxima were detected using rectangular dilation
with a neighbourhood approximately 0.9 expected radii wide. Peaks had to
reach at least 0.45 expected radii in distance. Connected peak plateaus
were represented by an actual plateau pixel nearest their centroid,
rather than assuming that a centroid necessarily lies on a peak.

Peaks were considered in decreasing distance order. A proposed marker
had to be at least 1.15 times the larger radius from every existing
Hough marker and previously accepted additional marker. It also had
to lie at least 0.6 expected radii from the crop edge. These checks
reduce duplicate markers and proposals truncated by local windows, but
can reject a distance peak displaced towards a neighbouring ball.

Markers used distinct positive integer labels in a 32-bit image. Hough
and additional foreground markers were disks of radius three pixels.
Sure background was outside a $7\times7$ elliptical dilation of the
foreground or outside the table polygon; the intervening collar remained
unknown. Watershed was applied to the original BGR crop after
$3\times3$ Gaussian smoothing with $\sigma=0.8$. It was not applied
directly to the floating-point distance transform. Separate seeded and
hybrid label maps were saved for comparison.

For each additional marker, its labelled basin was intersected with the
foreground mask. Acceptance required:
\begin{itemize}
\item Basin pixel area between 0.25 and 1.8 times the expected disk area.
\item Circularity $4\pi A/P^2\geq0.35$, measured from the largest external
contour's area $A$ and perimeter $P$.
\item Largest-contour bounding-box aspect ratio at most 2.2.
\item At least 35\% foreground support inside an expected-radius disk.
\item No duplicate of an existing or newly appended circle, and passage
through the same pocket filter used for baseline detections.
\end{itemize}

An accepted circle would use the distance peak as its centre and the
local median Hough radius as its radius. A visible watershed basin's
centroid was not substituted for an existing Hough centre, since partial
occlusion can shift a visible-region centroid away from the physical
ball centre. Existing circles and radii were preserved. Consequently,
this prototype could add missed candidates but could not correct an
existing inaccurate centre or delete a baseline false positive.

\subsubsection{Detection and recognition results}

% Generated by experiments/writeRecognitionReport.py from saved CSVs.
\begin{table}[htbp]
\centering
\small
\caption{Current-run detection comparison on Table1--6. The 90 annotations and 15-table-pixel matching tolerance are identical for all methods. Error is mean matched centre distance in table pixels.}
\label{tab:watershed-summary}
\begin{tabular}{lrrrrrrr}
\toprule
Method & TP & FP & FN & Precision & Recall & F1 & Error \\
\midrule
Hough baseline & 86 & 2 & 4 & 97.73\% & 95.56\% & 96.63\% & 4.91 \\
Hough-seeded watershed & 86 & 2 & 4 & 97.73\% & 95.56\% & 96.63\% & 4.91 \\
Hybrid watershed & 86 & 2 & 4 & 97.73\% & 95.56\% & 96.63\% & 4.91 \\
\bottomrule
\end{tabular}
\end{table}
\FloatBarrier

% Generated by experiments/writeRecognitionReport.py from saved CSVs.
\begin{table}[htbp]
\centering
\small
\caption{Distance-peak outcomes across all seven images. Only peaks passing marker placement proceed to watershed-basin validation.}
\label{tab:watershed-peaks}
\begin{tabular}{lr}
\toprule
Outcome & Count \\
\midrule
basin\_area & 1 \\
basin\_circularity & 3 \\
hough\_duplicate & 40 \\
near\_crop\_edge & 8 \\
\bottomrule
\end{tabular}
\end{table}
\FloatBarrier


The prototype did not recover any real-image detections. Across all
seven images it produced 52 eligible local distance maxima. Forty were
rejected as duplicates of Hough markers and eight lay too close to a
window edge. Four additional markers proceeded to basin validation:
three were rejected for insufficient circularity and one for insufficient
area. None was accepted. The resulting hybrid circle lists therefore
equalled the baseline circle lists.

On Table1--Table6, all methods consequently had 86 true positives,
two false positives and four false negatives. Precision was $97.73\%$,
recall $95.56\%$ and F1 $96.63\%$. Mean matched-centre error was
4.91 table pixels. No annotated ball was recovered or lost. The
Hough-seeded control was expected to retain baseline detection scores
by design; the hybrid's unchanged score was the empirical negative result.
AllBalls retained 16 matched detections with no false detections.

% Generated by experiments/writeRecognitionReport.py from saved CSVs.
\begin{table}[htbp]
\centering
\small
\caption{Per-image baseline and hybrid counts. Recognition is fixed to the tuned K=6 classifier plus selected solid-only CCA. AllBalls is the reference image.}
\label{tab:watershed-images}
\begin{tabular}{lrrrrrrr}
\toprule
Image & Base TP & FP & FN & Hybrid TP & FP & FN & Correct no. \\
\midrule
AllBalls.jpeg & 16 & 0 & 0 & 16 & 0 & 0 & 16 \\
Table1.jpeg & 14 & 1 & 2 & 14 & 1 & 2 & 10 \\
Table2.jpeg & 16 & 1 & 0 & 16 & 1 & 0 & 15 \\
Table3.jpeg & 15 & 0 & 1 & 15 & 0 & 1 & 14 \\
Table4.jpeg & 14 & 0 & 1 & 14 & 0 & 1 & 13 \\
Table5.jpeg & 14 & 0 & 0 & 14 & 0 & 0 & 12 \\
Table6.jpeg & 13 & 0 & 0 & 13 & 0 & 0 & 11 \\
\bottomrule
\end{tabular}
\end{table}
\FloatBarrier


Recognition was held fixed to configuration 2153 plus the selected
solid-only CCA correction from Section~\ref{sec:cca-search}. All methods
used the standard circular extraction mask, not the watershed basin.
This isolates detection proposals from an additional change to recognition
segmentation. With no accepted additional circles, correct type and
number counts remained 82/86 and 75/86 on matched Table1--Table6 balls,
and 16/16 on AllBalls. End-to-end number recall remained 75/90
($83.33\%$). This experiment did not test whether substituting watershed
masks for circular recognition masks would improve colour estimates.

\subsubsection{Why the additional markers failed}

The diagnostics showed a limitation of distance-transform seeding for
these particular masks. In Table1, closely packed balls, shadows and
foreground cleanup formed a broad connected foreground mass. Its
distance transform had large interior peaks corresponding to the group
shape, rather than a separate peak at every visible ball. The two
extracted peaks were both close to existing Hough detections. In
particular, there was no independent accepted marker for the missed
green or maroon ball. The green-cloth similarity and occluded boundaries
are plausible contributors, but the experiment did not separately
quantify their causal effects.

In the Table3 example, the missed striped blue ball touched the detected
black ball. Their combined foreground produced a distance maximum
between the visible objects, near the already detected black ball.
Duplicate suppression rejected it as an additional marker. A geometric
peak in the merged silhouette therefore did not reliably identify the
missing ball's centre. All four missed annotations lay inside processed
windows, so simply enlarging the processing coverage would not resolve
this observed marker problem.

\begin{figure}[htbp]
\centering
\includegraphics[width=\linewidth]{watershed_table1.png}
\caption{Table1 crowded-region diagnostics. The foreground mask merges
the packed group, and the distance transform does not provide one
independent marker per ball. Hough-seeded and hybrid watershed use the
same surviving markers in this case. The final panel shows rejected
distance peaks in red; no new circles were accepted. Watershed region
colours are arbitrary labels and white lines are watershed boundaries.}
\label{fig:watershed-table1}
\end{figure}
\FloatBarrier

\begin{figure}[htbp]
\centering
\includegraphics[width=0.78\linewidth]{watershed_table3.png}
\caption{A Table3 window containing the missed striped blue ball next
to a detected black ball. The merged foreground's distance peak does
not provide a sufficiently independent new marker. The seeded label
map consequently does not separate the missed ball as a new object.}
\label{fig:watershed-table3}
\end{figure}
\FloatBarrier

The implementation was also checked on a synthetic scene containing
three overlapping, differently coloured circles on a green background,
with only two Hough circles supplied. The hybrid recovered one additional
circle within three pixels of the known synthetic centre, preserved
the two supplied circles exactly, and retained the expected radius.
Additional checks covered empty or isolated detections, invalid image
types and insufficient cloth samples. This confirms that the proposal
path can function under controlled conditions; it does not establish
reliability on the real crowded scenes.

\subsubsection{Interpretation, outputs and reproducibility}

The measured conclusion is that this fixed watershed prototype did not
improve detection. It does not establish that every watershed-based
method would fail. The limiting factors were the provisional foreground
representation, the availability of independent markers, and the
conservative acceptance policy. Specular highlights, stripes and
neighbouring balls also create internal colour boundaries that need
not coincide with object boundaries. Further work could investigate
additional edge-based markers or fitting visible circular arcs, but
those alternatives were not implemented in this comparison. Acceptance
thresholds were not relaxed after inspecting the annotations in an
attempt to force a positive result.

There are no pixel-level ground-truth masks in the provided data.
Segmentation quality was therefore inspected visually; the experiment
reports centre-detection and recognition scores rather than mask IoU
or boundary accuracy. Watershed cannot directly reconstruct the hidden
part of a ball, and a segmented visible region is not automatically
an unbiased estimate of its complete circular boundary.

The experiment runs from the repository root with:
\begin{quote}
\code{./build-local/BallExperiments --watershed}
\end{quote}
An optional output-directory argument preserves separate runs. The
default directory is \code{output/experiments/watershed_detection/}.
The source options are in \code{experiments/watershedDetection.h}.
Saved files include pooled/per-image metrics, the raw-versus-filtered
ground-truth audit, local-window metadata, every peak's acceptance or
rejection reason, matched recognition predictions, and recovered/lost
annotation identities. Parameter values are exported to CSV and the
prototype procedure is described in a companion Markdown file.
Each image has one full detection comparison and one diagnostic sheet
per crowded window, rather than individual ball subdirectories.

Report snapshots are refreshed using
\code{python3 experiments/writeWatershedReport.py}. The surrounding
discussion contains run-specific observations and must be reviewed if
the implementation or results change. The main detection and recognition
pipelines remain unchanged.
